\section{Dificuldades Encontradas}

Durante o desenvolvimento do projeto, uma das principais dificuldades 
encontradas esteve relacionada à gestão de dependências do ecossistema 
\textit{Qiskit}, particularmente na integração entre os módulos 
\textit{qiskit}, \textit{qiskit-nature} e \textit{qiskit-algorithms}. Por se 
tratar de um framework em constante evolução, com atualizações frequentes e 
mudanças estruturais significativas (especialmente após a consolidação da 
versão 1.0), observou-se um desencontro de versionamento interno entre 
bibliotecas interdependentes.
Especificamente, ao utilizar rotinas do \textit{Qiskit Nature} para a 
obtenção da energia exata via \textit{Full Configuration Interaction} (FCI), 
ocorreu uma falha de importação associada à classe \texttt{BaseEstimator}, 
removida das versões mais recentes do módulo \textit{qiskit.primitives}. 
Embora o projeto não utilizasse diretamente funcionalidades de estados 
excitados ou do método QEOM, o erro era acionado por dependências internas 
do pacote, evidenciando como alterações estruturais em camadas não 
utilizadas diretamente podem impactar o funcionamento global do sistema.

Esse cenário ilustra uma característica recorrente em áreas de fronteira 
tecnológica: a rápida evolução dos frameworks pode introduzir instabilidades 
temporárias de compatibilidade entre módulos. A resolução adotada consistiu 
em reformular a estratégia de obtenção da energia de referência. Em vez de 
empregar o \textit{GroundStateEigensolver} do \textit{Qiskit Nature}, optou-se por realizar a diagonalização exata diretamente sobre o Hamiltoniano 
mapeado em qubits, utilizando o \textit{NumPyMinimumEigensolver}. 
Do ponto de vista físico e matemático, essa abordagem é equivalente ao 
cálculo de FCI, uma vez que corresponde à obtenção do menor autovalor do 
operador Hamiltoniano no espaço completo de Fock mapeado. Além de contornar 
o problema de dependência interna, essa solução tornou o pipeline 
computacional mais direto, reduzindo camadas intermediárias e aumentando a 
robustez do código frente a futuras atualizações do framework.

Assim, a principal dificuldade encontrada não esteve associada aos 
fundamentos físicos ou ao método variacional em si, mas sim à engenharia de 
software em um ambiente de desenvolvimento ativo e em transformação contínua,
aspecto característico da computação quântica contemporânea.